% Options for packages loaded elsewhere
\PassOptionsToPackage{unicode}{hyperref}
\PassOptionsToPackage{hyphens}{url}
%
\documentclass[
]{article}
\usepackage{amsmath,amssymb}
\usepackage{lmodern}
\usepackage{iftex}
\ifPDFTeX
  \usepackage[T1]{fontenc}
  \usepackage[utf8]{inputenc}
  \usepackage{textcomp} % provide euro and other symbols
\else % if luatex or xetex
  \usepackage{unicode-math}
  \defaultfontfeatures{Scale=MatchLowercase}
  \defaultfontfeatures[\rmfamily]{Ligatures=TeX,Scale=1}
\fi
% Use upquote if available, for straight quotes in verbatim environments
\IfFileExists{upquote.sty}{\usepackage{upquote}}{}
\IfFileExists{microtype.sty}{% use microtype if available
  \usepackage[]{microtype}
  \UseMicrotypeSet[protrusion]{basicmath} % disable protrusion for tt fonts
}{}
\makeatletter
\@ifundefined{KOMAClassName}{% if non-KOMA class
  \IfFileExists{parskip.sty}{%
    \usepackage{parskip}
  }{% else
    \setlength{\parindent}{0pt}
    \setlength{\parskip}{6pt plus 2pt minus 1pt}}
}{% if KOMA class
  \KOMAoptions{parskip=half}}
\makeatother
\usepackage{xcolor}
\usepackage{longtable,booktabs,array}
\usepackage{calc} % for calculating minipage widths
% Correct order of tables after \paragraph or \subparagraph
\usepackage{etoolbox}
\makeatletter
\patchcmd\longtable{\par}{\if@noskipsec\mbox{}\fi\par}{}{}
\makeatother
% Allow footnotes in longtable head/foot
\IfFileExists{footnotehyper.sty}{\usepackage{footnotehyper}}{\usepackage{footnote}}
\makesavenoteenv{longtable}
\setlength{\emergencystretch}{3em} % prevent overfull lines
\providecommand{\tightlist}{%
  \setlength{\itemsep}{0pt}\setlength{\parskip}{0pt}}
\setcounter{secnumdepth}{-\maxdimen} % remove section numbering
\ifLuaTeX
  \usepackage{selnolig}  % disable illegal ligatures
\fi
\IfFileExists{bookmark.sty}{\usepackage{bookmark}}{\usepackage{hyperref}}
\IfFileExists{xurl.sty}{\usepackage{xurl}}{} % add URL line breaks if available
\urlstyle{same} % disable monospaced font for URLs
\hypersetup{
  hidelinks,
  pdfcreator={LaTeX via pandoc}}

\author{}
\date{}

\begin{document}

\hypertarget{question-1}{%
\subsubsection{Question 1}\label{question-1}}

\hypertarget{a}{%
\paragraph{(a)}\label{a}}

truth table:

\begin{longtable}[]{@{}llllllll@{}}
\toprule()
\(r\) & \(p\) & \(q\) & \(\neg q\) & \(p\vee \neg q\) &
\(\color{purple} r\Rightarrow (p\vee \neg q)\) & \(q\wedge r\) &
\(\color{orange}(q\wedge r)\Rightarrow p\) \\
\midrule()
\endhead
T & T & T & F & T & T & T & T \\
T & T & F & T & T & T & F & T \\
T & F & T & F & F & F & T & F \\
F & T & T & F & T & T & F & T \\
T & F & F & T & T & T & F & T \\
F & T & F & T & T & T & F & T \\
F & F & T & F & F & T & F & T \\
T & T & T & F & T & T & T & T \\
\bottomrule()
\end{longtable}

By the 6th and 8th columns, which are exactly the same, we can conclude
that
\(\color{purple} r\Rightarrow (p\vee \neg q)\color{green}\equiv \color{orange}(q\wedge r)\Rightarrow p\)

\hypertarget{b}{%
\paragraph{(b)}\label{b}}

truth table:

\begin{longtable}[]{@{}lllllll@{}}
\toprule()
\(r\) & \(p\) & \(q\) & \(p\Rightarrow q\) &
\(\color{purple}(p\Rightarrow q)\Rightarrow r\) & \(q\Rightarrow r\) &
\(\color{orange}p\Rightarrow(q\Rightarrow r)\) \\
\midrule()
\endhead
T & T & T & T & T & T & T \\
T & T & F & F & T & T & T \\
T & F & T & T & T & T & T \\
F & T & T & T & F & F & F \\
T & F & F & T & T & T & T \\
F & T & F & F & T & T & T \\
F & F & T & T & F & F & T \\
F & F & F & T & F & T & T \\
\bottomrule()
\end{longtable}

By the last 2 rows, of which the 6th and 8th columns are different from
each other, we can conclude that
\(\color{purple}(p\Rightarrow q)\Rightarrow r\color{red} \not\equiv \color{orange}p\Rightarrow(q\Rightarrow r)\)

\hypertarget{question-2}{%
\subsubsection{Question 2}\label{question-2}}

The problem can be modeled as a constrained counting problem.

$$\left\{\begin{array}{l}10x_1+10x_2+10x_3=300\\
200\geqslant 10x_1,10x_2,10x_3\geqslant 50\\
\end{array}\right.\qquad \left\{\begin{array}{l}x_1+x_2+x_3=30\\
20\geqslant x_1,x_2,x_3\geqslant 5\end{array}\right.$$

let $y_i=x_i-4,i\in\{1,2,3\}$, then 

$$\left\{\begin{array}{l}y_1+y_2+y_3=18\\
16\geqslant y_1,y_2,y_3\geqslant 1\end{array}\right.$$

the equations can be considered as dividing $18$ same balls into $3$ non-empty distinct boxes, each box contains at most $16$ balls. 

Note that since the box cannot be empty, condition $16\geqslant y_1,y_2,y_3$ is abundant. Therefore the number of possible triples $(a_{\text{Alexander}},a_{\text{Kik}},a_{\text{Petros}})$ is

$$\displaystyle \binom{18-1}{3}=680$$.

\hypertarget{question-3}{%
\subsubsection{Question 3}\label{question-3}}

\begin{quote}
If \(x^2=y^2+z^2\), then \(xyz\) is even.
\end{quote}

Contrapositive of this implication:

\begin{quote}
If \(xyz\) is NOT even (i.e. is odd), then \(x^2\neq y^2+z^2\)
\end{quote}

Since the contrapositive statement is true if and only if the original
statement is true, vice versa.

Hence we\textquotesingle ll prove the original statement by proving its
contrapositive statement.

If \(xyz\) is odd, then \(x,y,z\) must be all odd. Hence \(x^2,y^2,z^2\)
are odd, which leads to that \(y^2+z^2\) is even but \(x^2\) is odd.
Therefore \(x^2\neq y^2+z^2\). Proved.

So we can conclude that if \(x^2=y^2+z^2\), then \(xyz\) is even.

\hypertarget{question-4}{%
\subsubsection{Question 4}\label{question-4}}

Given \(x\) is rational, suppose that \(x^2=x+1\). WLOG, let
\(x=\dfrac{p}{q}\), where \(p\in \mathbb{Z}\) and
\(q\in \mathbb{Z}-\{0\}\) are coprime. Hence
\(\dfrac{p^2}{q^2}=\dfrac{p}{q}+1\Leftrightarrow p^2=q(p+q)\).

By the fundamental theorem of arithmetic, the left hand side and right
hand side have exactly the same prime decompositions with exactly the
same multiplicities, hence \(\exists\) prime \(h\neq 1\), \(h\mid q\)
such that \(h\mid p\), then
\(1<h\mid \gcd (p,q)\Rightarrow \gcd(p,q) >1\).

This contradicts to the assumption that \(p,q\) are coprime. Therefore
there must be \(x^2\neq x+1\). Proved.

\hypertarget{question-5}{%
\subsubsection{Question 5}\label{question-5}}

There are 3 o\textquotesingle s, 2 t\textquotesingle s, 2
r\textquotesingle s in total, and one u, e, p, i, c each letter. Total
length 12.

all anagrams: \(\dfrac{12!}{3!2!2!}=19958400\)

anagrams without consecutive o\textquotesingle s:
\(\displaystyle \dfrac{(12-3)!}{2!2!}\times \binom{10}{3}=10886400\)

anagrams with at least 2 consecutive o\textquotesingle s:
\(19958400-10886400=9072000\)

\end{document}
